%% kapa: exam=1 round=23 year=2012 period=2 form=A subject=appraisal_law subject_ko=부동산관계법규 seq=17 number=57
%% source: Q-Net 감정평가사 기출문제 (cst00309) artlSeq=5052580 file=raw/1008/5052580/2012년 제23회 감정평가사 제1차시험 기출문제(2교시 A형).pdf
%% conversion: kordoc 4.24.51 -> tools/gen1.py (auto); review=proofread-sonnet; flags=none
\begin{question}{57}
부동산 가격공시 및 감정평가에 관한 법률의 목적과 관련하여 ( )의 용어에 관한 설명으로 옳지 않은 것은?
\begin{framedtext}
이 법은 토지, ( ㉠ ) 등 부동산의 ( ㉡ )을(를) 공시하여 부동산 가격산정의 기준이 되게 하고, 토지ㆍ건물ㆍ동산 등의 ( ㉢ )에 관한 사항을 정함으로써 이의 적정한 가격형성을 도모하며, 나아가 국토의 효율적인 이용과 국민경제의 발전에 이바지하게 함을 목적으로 한다. \par
\end{framedtext}
\begin{choices}
  \item ㉠에 포함되는 공동주택은 아파트, 연립주택, 다세대주택, 다가구주택으로 구분된다.
  \item ㉡은 통상적인 시장에서 정상적인 거래가 이루어지는 경우 성립될 가능성이 가장 높다고 인정되는 가격을 말한다.
  \item 표준지의 ㉡을 평가함에 있어 건물 그 밖의 정착물이 있는 때에는 당해 정착물이 존재하지 아니하는 것으로 보고 평가하여야 한다.
  \item 국토해양부장관이 표준지의 ㉡을 조사ㆍ평가하고자 할 때에는 둘 이상의 감정평가업자에게 이를 의뢰하여야 한다.
  \item ㉢은 토지등의 경제적 가치를 판정하여 그 결과를 가액으로 표시하는 것을 말한다.
\end{choices}
\end{question}
\answer{1}
